\section{Engineering reporting contract}\label{sec:reporting}

The empirical article must be interpretable as an AI-systems result, not only as a representation proposal. We therefore predeclare a reporting contract that couples capability, transfer validity and resource cost. The experimental unit is the task or evaluation item; repeated seeds are nested replicates. Principal comparisons freeze the base checkpoint, tokenizer, task set, held-out domains/families, context window, tool permissions, evaluator, budget schedule and seed schedule. Training-selection comparisons additionally freeze the downstream optimizer and training recipe; inference comparisons freeze the model weights.

\subsection{Metrics that answer the engineering questions}

For transfer discrimination, we report held-out ROC-AUC/PR-AUC, log loss or Brier score, calibration, Q2 true-accept and Q3 false-accept against the strongest semantic/skill control. The load-bearing quantity is incremental held-out value from witnessed structure, not raw accuracy of a structural classifier.

For training efficiency, let $T_m(q^*)$ be the minimum cumulative training-token exposure required by method $m$ to reach a preregistered structural-OOD capability target $q^*$. The headline token reduction against parent $p$ is
\[
\Delta_T(q^*)=1-\frac{T_{\rm RAKL}(q^*)}{T_p(q^*)}.
\]
We report the same target-based comparison for examples and, where measurable, training compute, together with a normalized area under the learning curve over a frozen token range. Structural-OOD is split into tasks with known structure but a new domain and the harder boundary in which whole structure families are held out; ordinary in-domain performance is retained as a safety check.

For inference efficiency, a \emph{valid solve} requires task correctness plus the registered transfer and verification gates. Billable inference tokens sum provider-reported input and output tokens across retrieval planning, adaptation, criticism and verification calls. For stratum $s$,
\[
\operatorname{TPVS}(s)=
\frac{\sum_{i\in s}\text{billable inference tokens}_i}
{\sum_{i\in s}\mathbf 1[\text{valid solve}_i]}.
\]
Zero valid solves therefore produce an infinite rather than missing cost-per-success value. We additionally report provider cost, tool calls, retrieval operations, median and p95 latency, and verification overhead.

The amortization claim uses cumulative total cost rather than online tokens alone:
\[
\begin{aligned}
C_{\rm total}(n)={}&C_{\rm induction}+C_{\rm training\ delta}+C_{\rm storage}\\
&+C_{\rm retrieval}(n)+C_{\rm adaptation}(n)+C_{\rm reasoning}(n)\\
&+C_{\rm tools}(n)+C_{\rm verification}(n).
\end{aligned}
\]
The empirical break-even count is the smallest registered reuse count for which RAKL is no more costly than the comparator while the capability target and invalid-transfer ceiling remain satisfied. A workload with no break-even in the registered horizon is reported as such.

Table~\ref{tab:transfer-metrics} summarizes the portable quantities an implementer can instrument on their own cross-domain reuse system. They are reported as a profile, never a single score, and describe transfer discipline and cost, not raw accuracy.

\begin{table}[t]
\footnotesize
\begin{tabular}{@{}p{0.24\textwidth} p{0.42\textwidth} p{0.28\textwidth}@{}}
\hline
\textbf{Metric} & \textbf{Definition} & \textbf{What it tells an engineer} \\
\hline
Invalid-transfer false-accept rate & verifier-invalid transfers the gate licensed, over all invalid transfers & whether the gate is truly fail-closed; the load-bearing safety number (target $0$) \\
Abstention coverage & cases resolved to \texttt{CANNOT\_CHECK}, over all cases & how often the gate honestly declines instead of guessing \\
Paired reduction vs.\ strongest control & held-out paired binary-Brier (or accuracy) reduction of the full contract over the best non-structural control, with bootstrap CI & whether witnessed structure adds value beyond the strongest baseline, not just beyond lexical similarity \\
Break-even count & smallest reuse count at which the structural layer is no costlier than the comparator; ``no finite break-even'' when reuse never pays & when reuse becomes economical for a workload \\
Tokens per valid solve (TPVS) & billable inference tokens over count of solves that also pass the transfer/verification gates ($\infty$ at zero valid solves) & cost per \emph{trustworthy} success, not per raw answer \\
\hline
\end{tabular}
\caption{Transfer-governance metrics an implementer can instrument. Reported as a profile; each measures transfer discipline or cost rather than accuracy alone.}
\label{tab:transfer-metrics}
\end{table}

\subsection{Main display logic}

The empirical manuscript is planned around six claims rather than six data tables. Figure~1 explains the shared structural object and matched training/inference design without making a performance claim. Figure~2 tests whether witnessed structure predicts valid transfer beyond semantic/skill controls, with Q2 acceptance and Q3 false acceptance shown together. Figure~3 reports structural-OOD performance versus cumulative training tokens and tokens/compute to a frozen capability target. Figure~4 reports the frozen-model valid-success versus inference-cost frontier, including tokens per valid solve and Q2/Q3 safety. Figure~5 reports cumulative total cost and the empirical break-even region after induction, retrieval, adaptation and verification are charged. Figure~6 reports directionality, QoI, boundary, evidence and verification ablations as effect sizes rather than a decorative feature inventory.

All quantitative displays use editable vector text, final-size uncertainty definitions and task-level source data. They contain no arrows pointing to data points, no opaque text boxes over data and no point-by-point prose annotations. Panel labels remain outside the data region. Method names are carried by fixed-axis labels or a shared legend rather than text placed over marks. Main figures show the shortest evidence chain needed for the central claim; per-model, per-family, cache/reuse, token-decomposition, calibration and sensitivity analyses are reserved for Extended Data unless they become load bearing after outcomes are observed.

\subsection{Representation ablations and shared-substrate criterion}

The structural object is only justified if its fields are load bearing. We therefore ablate directionality, QoI, target boundaries, evidence identity, explicit non-preserved properties and post-transfer verification, and compare against replacing the witness with semantic similarity or a domain-specific skill label. Primary ablation outcomes are Q2 acceptance, Q3 false acceptance, valid task success, tokens per valid solve and break-even count.

The strongest shared-substrate claim requires the same structural identity/object type to be used in both data selection and inference reuse. If training and inference require unrelated representations, the evidence supports two narrower engineering mechanisms rather than one unifying substrate. If the representation improves validity but never reaches a realistic total-cost break-even, the paper must describe it as a reliability mechanism rather than an efficiency mechanism.

The corresponding shared-substrate claim is falsified when an executed preregistered necessary condition fails. Missing confirmatory annotation is instead missing evidence: it leaves that claim unlicensed without converting an unrun test into a negative result.

The full machine-readable metric and figure contract is frozen before the model-level experiments are run.
